\chapter{Appendix: Algebra, Functions, and Proof}

Every limit, derivative, and integral in this book relies on older material:
rearranging algebraic expressions, the laws of exponents, inequalities and
absolute values, basic facts about functions, and the habit of proving that a
step is true. This appendix collects that material in one place, as a reference
for whenever a calculation in the main text seems puzzling.

For each rule we ask \emph{why} it holds, draw a picture where one helps, and
point out the places where students commonly slip. None of it is arbitrary: the
exponent laws follow once we decide what $a^n$ should mean; the rules for
inequalities follow from one fact about the number line; and a proof is an
explanation with no gaps.

\begin{focusbox}
Three things to carry away from this appendix:
\begin{enumerate}
  \item Algebraic identities (factoring, exponent laws, the binomial pattern)
        are not arbitrary formulas to memorize --- each is forced by a single
        defining idea, and once you see the idea you can reconstruct the rule.
  \item Inequalities and absolute values are the native language of calculus
        (limits, error bounds, convergence). The one rule everyone forgets ---
        multiplying by a negative flips the sign --- has a clean reason.
  \item A proof is an explanation with no gaps. Direct proof, proof by
        contradiction, and induction are three reliable shapes that explanation
        can take.
\end{enumerate}
\end{focusbox}

\section{Numbers, and the Ground We Stand On}
\label{ch40:numbers}

Before algebra there are the numbers algebra rearranges. We use a nested tower of
number systems, each built to fix a shortcoming of the last:
\[
  \N \subset \Z \subset \Q \subset \R \subset \C.
\]
Here $\subset$ means ``is contained in'': every natural number is an integer,
every integer is a rational number, and so on. Figure~\ref{ch40:fig-tower} draws
the tower as nested boxes. The naturals $\N = \set{1,2,3,\dots}$ count discrete items. The integers
$\Z = \set{\dots,-2,-1,0,1,2,\dots}$ add the ability to subtract freely. The
rationals $\Q$ --- ratios $p/q$ of integers with $q \neq 0$ --- add the ability to
divide. The reals $\R$ fill in the gaps between rationals so that the number line
has no holes (a number like $\sqrt 2$ or $\pi$ lives here but not in $\Q$). The
complex numbers $\C$ add a square root of $-1$. Calculus lives almost entirely in
$\R$, and the ``no holes'' property of $\R$ --- its \emph{completeness} (the
guarantee that the number line has no missing gaps or punctures, so every sequence
of numbers squeezing closer together lands on an existing real number) ---
is what makes limits behave; that deeper story has its own dedicated appendix.

\begin{figure}[htbp]
\centering
\begin{tikzpicture}[font=\small]
  \draw[rounded corners, thick, fill=structurepurple!8]  (0,0)     rectangle (13,4);
  \draw[rounded corners, thick, fill=sourceblue!10]      (0.3,0.4) rectangle (10.4,3.6);
  \draw[rounded corners, thick, fill=targetgreen!12]     (0.6,0.8) rectangle (7.8,3.2);
  \draw[rounded corners, thick, fill=quotientorange!15]  (0.9,1.2) rectangle (5.2,2.8);
  \draw[rounded corners, thick, fill=white]              (1.15,1.5) rectangle (3.45,2.5);
  \node[align=center] at (2.3,2.0)   {$\N$: $1,2,3,\dots$};
  \node[align=center] at (4.3,2.0)   {$\Z$\\ $-3,\ 0$};
  \node[align=center] at (6.5,2.0)   {$\Q$\\ $\tfrac12,\ -0.75$};
  \node[align=center] at (9.1,2.0)   {$\R$\\ $\sqrt2,\ \pi$};
  \node[align=center] at (11.7,2.0)  {$\C$\\ $i,\ 2+3i$};
\end{tikzpicture}
\caption{The number systems as nested boxes. Each box contains every box inside
it, and the examples written in a box belong to that system but not to the
smaller ones: $-3$ is an integer but not a natural number, $\sqrt2$ is real but
not rational.}
\label{ch40:fig-tower}
\end{figure}

\begin{historybox}
The language of algebra, number systems, and formal proof was not handed down
fully formed; it was worked out over about two thousand years.

Ancient and medieval mathematicians worked without modern symbolic notation. In the
third century CE, the Greek mathematician Diophantus of Alexandria solved algebraic
equations using words and abbreviations, but he dismissed an equation like $4x + 20 = 4$
as ``absurd'' because it demanded a negative answer. Around 825~CE, the Persian
scholar Muhammad ibn Musa al-Khwarizmi wrote \emph{al-Kitab al-mukhtasar fi hisab
al-jabr wa'l-muqabala} (``The Compendious Book on Calculation by Restoration and
Balancing''). The word \emph{al-jabr} (restoration, meaning moving a subtracted quantity
to the opposite side of an equation to make it positive) gave us our word
\emph{algebra}. Al-Khwarizmi and his contemporaries solved quadratics by literal
geometric cut-and-paste dissections of physical squares, but they avoided negative
quantities entirely because a square cannot have a negative side length.

For centuries, negative numbers remained mathematical outcasts. As late as the
seventeenth century, Ren\'e Descartes termed negative polynomial roots
``false roots'' (\emph{racines fausses}), and Blaise Pascal dismissed subtracting
$4$ from $0$, arguing that ``nothing can be less than nothing'' in nature. It was not
until John Wallis popularized the geometric number line in 1685 --- interpreting
positive and negative numbers as directed steps right and left from an
origin --- that negative quantities gained an intuitive home.

Meanwhile, algebra moved from words to symbols. In 1591,
Fran\c{c}ois Vi\`ete introduced the convention of using letters for both
unknowns and known parameters, turning algebra from a catalog of recipes written
in words into a symbolic language. But the symbols ran ahead of the logic.
During the early development of calculus, mathematicians manipulated
infinitesimals (quantities imagined to be smaller than every positive number yet
not zero), cancelled vanishing factors, and added up divergent series (infinite
sums with no finite total) by formal rules alone, prompting Bishop George Berkeley in \emph{The Analyst} (1734) to criticize
dividing by increments $\dd x \neq 0$ that were subsequently discarded as zero.
Manipulating formulas without clear definitions
frequently led to contradictions, such as bogus derivations that $1 = 2$ or that
$1 - 1 + 1 - 1 + \cdots = 1/2$.

In the nineteenth century, Augustin-Louis Cauchy, Karl
Weierstrass, and Richard Dedekind replaced the vague geometric intuitions of calculus
with explicit definitions built from inequalities, the completeness of the real
numbers, and step-by-step proof. Mathematical induction, used by Pascal
and made formal by Giuseppe Peano, gave a way to prove a statement for infinitely
many whole numbers at once (Section~\ref{ch40:proof}). Symbolic algebra and
deductive proof together gave science a precise language for deriving and
checking predictions about planetary orbits, quantum wavefunctions, and fluid
dynamics.

The question of what symbol manipulation can and cannot establish is still
studied. In 1931, Kurt G\"odel proved his incompleteness theorems: any
consistent, mechanically checkable system of axioms (basic rules assumed without
proof) strong enough for ordinary arithmetic contains true statements that cannot
be proved from those axioms. Today, proof assistants (software such as Lean and
Coq that checks every step of a proof down to the basic rules of logic) are being
used to re-check large bodies of mathematics.
\end{historybox}

\begin{intuitionbox}
Picture the real line as an infinitely long, perfectly smooth ruler. Every point
on it is a real number, and every real number sits at exactly one point. The
rationals are dense --- between any two of them there is another --- yet they leave
gaps you cannot see but can prove are there. The reals are what you get when you
plug every gap. Most of calculus is the art of saying precise things about points
sitting infinitely close together on this ruler.
\end{intuitionbox}

The arithmetic of $\R$ obeys a short list of \emph{field axioms}: basic rules,
assumed rather than proved, for how numbers add and multiply.
Addition and multiplication are \emph{commutative} (order does not matter:
$a+b=b+a$ and $ab=ba$) and \emph{associative} (grouping does not matter:
$(a+b)+c=a+(b+c)$, and likewise for products). The numbers $0$ and $1$ are
\emph{identities}: adding $0$ or multiplying by $1$ leaves a number unchanged.
Every number $a$ has an \emph{additive inverse} $-a$ (with $a+(-a)=0$), and every
nonzero number has a \emph{multiplicative inverse} $1/a$ (with $a\cdot\tfrac1a=1$).
Connecting the two operations, multiplication distributes over addition:
\[
  a(b+c) = ab + ac.
\]
That last law, the \emph{distributive law}, is the most used fact in algebra. Expanding brackets, factoring, multiplying polynomials, the
binomial theorem --- all of it is the distributive law applied again and again.

\begin{happeningbox}
When you ``expand $3(x+2)$ to get $3x+6$'' and when you ``factor $3x+6$ back to
$3(x+2)$'' you are doing the exact same equation read in two directions:
$a(b+c) = ab+ac$. Expanding goes left to right; factoring goes right to left.
This is why factoring feels harder --- running a rule \emph{backwards} means
searching for the inputs that would have produced what you see, and search is
always harder than evaluation.
\end{happeningbox}

\section{Algebraic Manipulation: Expanding and Factoring}
\label{ch40:factoring}

\subsection{Multiplying out}

To multiply two sums, distribute every term of the first across every term of the
second. For two \emph{binomials} (sums of two terms) this is the familiar pattern
(sometimes taught as ``FOIL'': first, outer, inner, last):
\[
  (a+b)(c+d) = ac + ad + bc + bd.
\]
Each term on the right is one term from the first factor times one term from the
second; we collect every such pairing. There is nothing special about two terms.
$(a+b)(c+d+e)$ produces six products, one for each (left term, right term) pair.

The identities worth knowing cold are the ones that appear constantly in
calculus, especially when rationalizing or simplifying difference quotients:
\begin{align}
  (a+b)^2 &= a^2 + 2ab + b^2, \label{ch40:square-sum}\\
  (a-b)^2 &= a^2 - 2ab + b^2, \\
  (a+b)(a-b) &= a^2 - b^2, \label{ch40:diff-squares}\\
  (a+b)^3 &= a^3 + 3a^2 b + 3ab^2 + b^3.
\end{align}
The difference of squares~\eqref{ch40:diff-squares} is especially useful: it
is the basis for multiplying by a \emph{conjugate} (the same two terms with the
sign between them flipped, such as $\sqrt{x+1}+1$ for $\sqrt{x+1}-1$), which you
meet the first time you compute a limit like $\lim_{x\to 0}(\sqrt{x+1}-1)/x$.

\begin{whybox}
Why is $(a+b)^2 = a^2 + 2ab + b^2$ and \emph{not} $a^2 + b^2$? Because squaring a
sum is not the same as summing the squares: the \emph{cross terms} (products that
mix $a$ with $b$) do not go away. Run the distributive law all the way out:
\[
  (a+b)^2 = (a+b)(a+b) = a\cdot a + a \cdot b + b \cdot a + b \cdot b
          = a^2 + 2ab + b^2.
\]
The middle $2ab$ is the two mixed pairings $ab$ and $ba$. Forgetting it is
probably the most common algebra error in calculus.
\end{whybox}

The cross term has a picture: Figure~\ref{ch40:fig-square} shows it as the two
rectangles you get when you enlarge a square.

\begin{figure}[htbp]
\centering
\begin{tikzpicture}[scale=1.1]
  % big square split into 4 regions
  \draw[fill=sourceblue!18] (0,0) rectangle (3,3);
  \draw[fill=quotientorange!25] (3,0) rectangle (4.2,3);
  \draw[fill=quotientorange!25] (0,3) rectangle (3,4.2);
  \draw[fill=targetgreen!22] (3,3) rectangle (4.2,4.2);
  \draw[thick] (0,0) rectangle (4.2,4.2);
  \draw[thick] (3,0) -- (3,4.2);
  \draw[thick] (0,3) -- (4.2,3);
  % labels of regions
  \node at (1.5,1.5) {$a^2$};
  \node at (3.6,1.5) {$ab$};
  \node at (1.5,3.6) {$ab$};
  \node at (3.6,3.6) {$b^2$};
  % side labels
  \node[below] at (1.5,-0.05) {$a$};
  \node[below] at (3.6,-0.05) {$b$};
  \node[left]  at (-0.05,1.5) {$a$};
  \node[left]  at (-0.05,3.6) {$b$};
\end{tikzpicture}
\caption{The square on side $a+b$ has area $(a+b)^2$. Cut it into four pieces:
the $a$-square, the $b$-square, and \emph{two} $a\times b$ rectangles. That is why
$(a+b)^2 = a^2 + 2ab + b^2$ --- the $2ab$ is the two orange rectangles, the part
people forget.}
\label{ch40:fig-square}
\end{figure}

\subsection{Factoring: running the machine backwards}

Factoring rewrites a sum as a product. It matters because a product is zero
exactly when one of its factors is zero (the \emph{zero-product property}), which
is how we solve equations, and because a factored expression often lets a
troublesome common factor cancel, which is the main way to compute limits of the
indeterminate form $0/0$.

The basic moves, in rough order of how often you need them:

\textbf{(1) Pull out a common factor.}
\[
  6x^3 - 9x^2 = 3x^2(2x - 3).
\]
Look for the largest factor shared by every term and lift it out front.

\textbf{(2) Difference of squares.}
\[
  x^2 - 25 = (x-5)(x+5), \qquad x^4 - 1 = (x^2-1)(x^2+1) = (x-1)(x+1)(x^2+1).
\]

\textbf{(3) Quadratic trinomials} (three-term expressions whose highest power is
$x^2$). To factor $x^2 + bx + c$, find two numbers that
\emph{multiply to $c$ and add to $b$}. For $x^2 + 7x + 12$ we want a pair with
product $12$ and sum $7$: that is $3$ and $4$, so
\[
  x^2 + 7x + 12 = (x+3)(x+4).
\]

\textbf{(4) Sum and difference of cubes.}
\[
  a^3 - b^3 = (a-b)(a^2 + ab + b^2), \qquad
  a^3 + b^3 = (a+b)(a^2 - ab + b^2).
\]

\begin{trapbox}
There is no real factoring of a \emph{sum} of squares: $a^2 + b^2$ does not factor
over the reals, that is, into factors whose coefficients are real numbers (you
need complex numbers). Students who have just factored $a^2 - b^2$
often write $a^2+b^2 = (a+b)(a-b)$ out of momentum --- but expanding that gives
$a^2 - b^2$, the wrong sign. The difference factors; the sum does not.
\end{trapbox}

\begin{example}[Factoring to resolve a $0/0$ limit]
\label{ch40:ex-factor-limit}
Consider $\dfrac{x^2 - 9}{x - 3}$ as $x \to 3$. Plugging in $x=3$ gives $0/0$,
which says nothing yet. Factor the top as a difference of squares:
\[
  \frac{x^2 - 9}{x-3} = \frac{(x-3)(x+3)}{x-3} = x + 3 \quad (x \neq 3).
\]
The offending factor $x-3$ cancels, and the simplified expression $x+3$ tends to
$6$ as $x \to 3$. The $0/0$ indeterminacy came from a single shared factor $(x-3)$,
which factoring exposed and cancelled.
\end{example}

\begin{trapbox}
\textbf{The danger of formal cancellation without domain vigilance.}
Algebraic formulas manipulate symbols blindly unless you track domain restrictions.
Consider the famous bogus ``proof'' that $1 = 2$:
Start with $a = b$. Multiply by $a$ to get $a^2 = ab$. Subtract $b^2$ from both sides:
$a^2 - b^2 = ab - b^2$. Factor both sides: $(a-b)(a+b) = b(a-b)$. Now cancel the
common factor $(a-b)$: $a+b = b$. Since $a=b$, this says $2b = b$, so dividing by $b$
gives $2 = 1$.

Where did the reasoning derail? At the cancellation step: because $a=b$, the factor
$a-b$ is identically $0$. Dividing both sides by zero destroyed the equation.
In calculus, when we write $\lim_{x\to 3}\frac{x^2-9}{x-3} = \lim_{x\to 3}(x+3)$,
the cancellation is valid precisely because the limit process considers values of
$x$ arbitrarily close to $3$, but \emph{never} $x = 3$. Hence $x - 3 \neq 0$ throughout
the approach.
\end{trapbox}

\begin{walkbox}
Try factoring $2x^2 + x - 6$. (Hint: look for two numbers multiplying to
$2\cdot(-6) = -12$ and adding to $1$; they are $4$ and $-3$. Split the middle term:
$2x^2 + 4x - 3x - 6 = 2x(x+2) - 3(x+2) = (2x-3)(x+2)$.) The technique of
``splitting the middle term'' turns a hard trinomial into two easy common-factor
pulls.
\end{walkbox}

\subsection{The quadratic formula, and completing the square}

Not every quadratic factors with nice integers. For \emph{any} quadratic
$ax^2 + bx + c = 0$ with $a \neq 0$, the solutions are
\[
  x = \frac{-b \pm \sqrt{b^2 - 4ac}}{2a}.
\]
The formula comes from \emph{completing the square}, a technique worth seeing in
full because the same move reappears in integration (trig substitution) and in
finding the vertex (lowest or highest point) of a parabola.

\begin{whybox}
Start from $ax^2 + bx + c = 0$. Divide by $a$ and move the constant:
\[
  x^2 + \tfrac{b}{a}x = -\tfrac{c}{a}.
\]
The left side is \emph{almost} a perfect square. The square $(x + \tfrac{b}{2a})^2$
expands to $x^2 + \tfrac{b}{a}x + \tfrac{b^2}{4a^2}$, so add $\tfrac{b^2}{4a^2}$ to
both sides to complete it:
\[
  \Bigl(x + \tfrac{b}{2a}\Bigr)^2 = \frac{b^2}{4a^2} - \frac{c}{a}
                                  = \frac{b^2 - 4ac}{4a^2}.
\]
Now take square roots (the square root of the denominator gives $\sqrt{4a^2} = 2\abs{a}$,
and absorbing the absolute value into the $\pm$ leaves the familiar $2a$) and solve for $x$:
\[
  x + \frac{b}{2a} = \pm\frac{\sqrt{b^2-4ac}}{2a}
  \;\Longrightarrow\;
  x = \frac{-b \pm \sqrt{b^2-4ac}}{2a}.
\]
The quadratic formula is completing the square done once, symbolically, so
you never have to repeat the algebra by hand.
Figure~\ref{ch40:fig-complete-square} shows the same step as a picture.
\end{whybox}

\begin{figure}[htbp]
\centering
\begin{tikzpicture}[scale=1.1]
  \draw[fill=sourceblue!20] (0,0) rectangle (3,3);
  \draw[fill=quotientorange!25] (3,0) rectangle (4.2,3);
  \draw[fill=quotientorange!25] (0,3) rectangle (3,4.2);
  \draw[fill=targetgreen!25, dashed, thick] (3,3) rectangle (4.2,4.2);
  \draw[thick] (0,0) rectangle (4.2,4.2);
  \draw[thick] (3,0) -- (3,3);
  \draw[thick] (0,3) -- (3,3);
  \node at (1.5,1.5) {$x^2$};
  \node at (3.6,1.5) {$\frac{b}{2}x$};
  \node at (1.5,3.6) {$\frac{b}{2}x$};
  \node[targetgreen!60!black, font=\small] at (3.6,3.6) {$(b/2)^2$};
  \node[below] at (1.5,-0.05) {$x$};
  \node[below] at (3.6,-0.05) {$b/2$};
  \node[left]  at (-0.05,1.5) {$x$};
  \node[left]  at (-0.05,3.6) {$b/2$};
  \draw[decorate,decoration={brace,amplitude=5pt,mirror},thick]
       (0,-0.5) -- (4.2,-0.5) node[midway,below=4pt]{total side length $x + b/2$};
\end{tikzpicture}
\caption{Completing the square seen as physical geometry. To make $x^2 + bx$ into a perfect square, divide the rectangular area $bx$ into two equal strips of width $b/2$ and place them along the right and top edges of the $x\times x$ square. The resulting $L$-shape is missing only the small top-right corner of area $(b/2)^2$ (green, dashed). Adding this missing piece completes a perfect square of total side $x + b/2$ and area $(x + b/2)^2$.}
\label{ch40:fig-complete-square}
\end{figure}

\begin{connectionbox}
Completing the square began as geometry. Four thousand years ago,
Babylonian scribes solved quadratic problems by cutting and rearranging physical areas.
When al-Khwarizmi solved $x^2 + 10x = 39$ in 825~CE, he drew a square of side $x$ and two
rectangles of width $5$, placed them along the sides of the square, and noted that the
entire shape was short of a full square by a small corner of area $5^2 = 25$. Adding $25$
to both sides completed the square: $(x+5)^2 = 39 + 25 = 64$, so
$x + 5 = 8$ and $x = 3$. (He accepted only the positive side length; the other
solution, $x+5=-8$, gives $x=-13$.) Figure~\ref{ch40:fig-complete-square} shows
this cut-and-paste step.
\end{connectionbox}

The quantity $b^2 - 4ac$ under the root is the \emph{discriminant} (it
discriminates, or tells apart, the possible kinds of solution). If it is
positive there are two distinct real roots; if zero, one repeated real root; if
negative, no real roots (the parabola never touches the $x$-axis), and the two
solutions are \emph{complex conjugates}: a pair $p+qi$ and $p-qi$, where
$i=\sqrt{-1}$. Figure~\ref{ch40:fig-discriminant} shows the three cases.

\begin{figure}[htbp]
\centering
\begin{tikzpicture}
\begin{axis}[
    name=two, width=4.9cm, height=4.6cm,
    axis lines=middle, xmin=-2.2, xmax=2.2, ymin=-1.6, ymax=3.6,
    xtick={-1,1}, ytick=\empty, domain=-2:2, samples=60,
    xlabel={$x$}, xlabel style={anchor=west},
    title={\small $b^2-4ac>0$\\ two roots}, title style={align=center},
    tick label style={font=\small},
]
  \addplot[curveblue, very thick] {x^2-1};
  \addplot[only marks, mark=*, mark size=1.8pt, tangentred] coordinates {(-1,0) (1,0)};
\end{axis}
\begin{axis}[
    name=one, at={(two.east)}, anchor=west, xshift=0.6cm,
    width=4.9cm, height=4.6cm,
    axis lines=middle, xmin=-2.2, xmax=2.2, ymin=-1.6, ymax=3.6,
    xtick=\empty, ytick=\empty, domain=-2:2, samples=60,
    xlabel={$x$}, xlabel style={anchor=west},
    title={\small $b^2-4ac=0$\\ one root}, title style={align=center},
]
  \addplot[curveblue, very thick] {x^2};
  \addplot[only marks, mark=*, mark size=1.8pt, tangentred] coordinates {(0,0)};
\end{axis}
\begin{axis}[
    at={(one.east)}, anchor=west, xshift=0.6cm,
    width=4.9cm, height=4.6cm,
    axis lines=middle, xmin=-2.2, xmax=2.2, ymin=-1.6, ymax=3.6,
    xtick=\empty, ytick=\empty, domain=-2:2, samples=60,
    xlabel={$x$}, xlabel style={anchor=west},
    title={\small $b^2-4ac<0$\\ no real roots}, title style={align=center},
]
  \addplot[curveblue, very thick] {x^2+1};
\end{axis}
\end{tikzpicture}
\caption{The sign of the discriminant decides how often the parabola
$y=ax^2+bx+c$ meets the $x$-axis. Left to right: $y=x^2-1$ crosses twice (at
$x=\pm1$), $y=x^2$ touches once, and $y=x^2+1$ never meets the axis.}
\label{ch40:fig-discriminant}
\end{figure}

\section{Exponents and Their Laws}
\label{ch40:exponents}

\subsection{Where the laws come from}

For a positive integer $n$, $a^n$ means ``multiply $a$ by itself $n$ times.'' From
this single definition every exponent law follows by counting factors. Take
$a^m \cdot a^n$: that is $m$ copies of $a$ times $n$ copies of $a$, which is $m+n$
copies. Hence
\[
  a^m \cdot a^n = a^{m+n}.
\]
There is nothing to memorize; you are counting how many $a$'s stand in the
line. The same counting gives the rest:
\begin{align*}
  a^m \cdot a^n &= a^{m+n}, & \frac{a^m}{a^n} &= a^{m-n}, \\
  (a^m)^n &= a^{mn}, & (ab)^n &= a^n b^n, \\
  \Bigl(\frac{a}{b}\Bigr)^n &= \frac{a^n}{b^n}. & &
\end{align*}

\begin{happeningbox}
Every exponent law stems from the fundamental identity $a^{m+n}=a^m a^n$,
which reflects how exponentiation turns addition (in the exponent) into
multiplication (of the values). $(a^m)^n = a^{mn}$ is repeated addition of
exponents: $n$ copies of $m$. This addition-becomes-multiplication structure is
what makes the exponential function $e^x$ so central in calculus: it is
the function that converts the additive world of inputs into the multiplicative
world of outputs, and, up to a constant multiple, it is the only function that is
its own derivative.
\end{happeningbox}

\subsection{Extending to zero, negatives, and fractions}

The definition ``multiply $a$ by itself $n$ times'' only makes sense for positive
integer $n$. So what could $a^0$, $a^{-3}$, or $a^{1/2}$ possibly mean? One
principle settles all three cases at once: \emph{we extend the definition in the only way that
keeps the law $a^m a^n = a^{m+n}$ true.} The laws are not re-proved for the new
cases; they are the very thing we insist on preserving, and that insistence forces
a unique answer.

\textbf{Zero exponent.} We need $a^0 \cdot a^n = a^{0+n} = a^n$. Dividing by $a^n$
(for $a \neq 0$) forces
\[
  a^0 = 1.
\]

\textbf{Negative exponent.} We need $a^{-n} \cdot a^n = a^{0} = 1$, so $a^{-n}$
must be the reciprocal:
\[
  a^{-n} = \frac{1}{a^n}.
\]

\textbf{Fractional exponent.} Take $a > 0$ and a positive integer $n$. We need
$\bigl(a^{1/n}\bigr)^n = a^{(1/n)\cdot n} = a^1 = a$, so $a^{1/n}$ must be the
positive number whose $n$th power is $a$ --- the $n$th root:
\[
  a^{1/n} = \sqrt[n]{a}, \qquad a^{m/n} = \bigl(\sqrt[n]{a}\bigr)^m = \sqrt[n]{a^m}
  \qquad (a > 0).
\]
The restriction $a > 0$ matters. For $a < 0$ an even root is not a real number, and
the exponent laws break: $\bigl((-1)^2\bigr)^{1/2} = 1$, while
$(-1)^{2\cdot\frac12} = -1$.

\begin{whybox}
Mathematics uses this move again and again.
The expression $a^{1/2}$ has \emph{no} prior meaning --- you cannot multiply $a$
by itself ``half a time.'' We \emph{define} it. But we do not define it by whim:
we choose the one value that lets the exponent law survive. That demand pins it
down uniquely to $\sqrt a$. Mathematics grows by exactly this method: take a rule
that holds in a small world, then extend the world in whatever way keeps the rule
intact. The same logic later defines $a^x$ for irrational $x$, $0! = 1$, and
$e^{i\theta} = \cos\theta + i\sin\theta$. Figure~\ref{ch40:fig-exp} shows the
result for base $2$.
\end{whybox}

\begin{figure}[htbp]
\centering
\begin{tikzpicture}
\begin{axis}[
    width=11.5cm, height=6.4cm,
    xlabel={$x$}, ylabel={$2^x$},
    xlabel style={at={(1,0)}, anchor=south west},
    ylabel style={at={(ticklabel* cs:1)}, anchor=south},
    domain=-3:3, samples=120,
    xmin=-3.2, xmax=3.3, ymin=-0.4, ymax=9.2, ytick={2,4,6,8},
    axis lines=middle,
    grid=major, grid style={gray, opacity=0.2},
    legend pos=north west,
]
  \addplot[curveblue, very thick] {2^x};
  \addlegendentry{$y = 2^x$}
  % mark special integer points
  \addplot[only marks, mark=*, mark size=1.6pt, tangentred]
     coordinates {(-2,0.25) (-1,0.5) (0,1) (1,2) (2,4) (3,8)};
  \node[anchor=north west, font=\small] at (axis cs:0.1,0.95) {$2^0=1$};
  \node[anchor=south east, font=\small] at (axis cs:-2,0.32) {$2^{-2}=1/4$};
\end{axis}
\end{tikzpicture}
\caption{The exponential $2^x$ knits the integer powers (red dots) into a smooth
curve. Negative exponents give reciprocals ($2^{-2}=\tfrac14$), the zero exponent
gives $1$, and fractional exponents fill in everything between. The values for
non-integer $x$ are precisely the ones that keep $2^{m+n}=2^m2^n$ holding for all
real $m,n$.}
\label{ch40:fig-exp}
\end{figure}

\begin{trapbox}
The exponent laws govern \emph{products with a shared base}, not sums and not
mixed bases. All of the following are wrong:
\[
  a^m + a^n \neq a^{m+n}, \qquad
  (a+b)^n \neq a^n + b^n, \qquad
  \sqrt{a^2+b^2} \neq a + b.
\]
The last is the same error as the missing cross term in $(a+b)^2$. There is no law
that simplifies $a^m + a^n$ in general --- the laws live on the multiplicative side
only. Whenever you are tempted to ``distribute'' an exponent across a sum, stop:
exponents distribute over products, never over sums.
\end{trapbox}

\begin{recoverybox}
\textbf{Pause and think:} Why is $(-2)^4 = 16$, but $-2^4 = -16$?
By standard order of operations, exponentiation binds tighter than negation. In
$(-2)^4$, the parentheses bundle the negative sign into the base:
$(-2)(-2)(-2)(-2) = +16$. In $-2^4$, the power $2^4 = 16$ is computed first, and
the minus sign is placed out front, giving $-16$. A missing set of parentheses
in a derivative computation is one of the most common ways to accidentally flip
a sign and derail an entire solution.
\end{recoverybox}

\section{Inequalities}
\label{ch40:inequalities}

Equations ask \emph{where} two things are equal; inequalities ask \emph{which is
bigger}. Calculus runs on inequalities far more than on equations: the definition
of a limit is an inequality ($\abs{f(x)-L}<\eps$), comparison tests for series are
inequalities, error bounds are inequalities.

The order relation $<$ on $\R$ obeys a handful of rules. Most are unsurprising, but
one is the famous trap.

\begin{focusbox}
\textbf{Rules for inequalities.} For real numbers,
\begin{itemize}
  \item \textbf{Add anything:} if $a < b$ then $a + c < b + c$ (for any $c$).
  \item \textbf{Multiply by a positive:} if $a<b$ and $c>0$ then $ac < bc$.
  \item \textbf{Multiply by a negative --- FLIP:} if $a<b$ and $c<0$ then
        $ac > bc$.
  \item \textbf{Transitivity:} if $a<b$ and $b<c$ then $a<c$.
  \item \textbf{Reciprocal of positives flips:} if $0<a<b$ then
        $\tfrac1a > \tfrac1b$.
\end{itemize}
\end{focusbox}

\begin{whybox}
Why does multiplying by a negative \emph{flip} the inequality? Picture the number
line. The statement $a<b$ means $a$ sits to the left of $b$. Multiplying every
number by $-1$ is a \emph{reflection of the line through $0$}: left and right swap.
So whatever was on the left ($a$) lands on the right, and the order reverses. More
algebraically: $a<b$ means $b-a>0$, i.e.\ $b-a$ is positive. Multiply a positive
number by a negative $c$ and you get a negative: $c(b-a)<0$, that is $cb-ca<0$,
i.e.\ $cb<ca$, i.e.\ $ca>cb$. The flip is forced.
Figure~\ref{ch40:fig-flip} shows the reflection.
\end{whybox}

\begin{figure}[htbp]
\centering
\begin{tikzpicture}[scale=1.0]
  % original line
  \draw[->,thick] (-4.6,0) -- (4.6,0);
  \foreach \x in {-4,...,4} \draw (\x,0.08)--(\x,-0.08);
  \fill[sourceblue] (1,0) circle (2.6pt) node[above=3pt]{$a=1$};
  \fill[targetgreen] (3,0) circle (2.6pt) node[above=3pt]{$b=3$};
  \node at (-5.4,0) {original};
  % reflected line below
  \draw[->,thick] (-4.6,-1.8) -- (4.6,-1.8);
  \foreach \x in {-4,...,4} \draw (\x,-1.72)--(\x,-1.88);
  \fill[sourceblue] (-1,-1.8) circle (2.6pt) node[below=3pt]{$-a=-1$};
  \fill[targetgreen] (-3,-1.8) circle (2.6pt) node[below=3pt]{$-b=-3$};
  \node at (-5.4,-1.8) {$\times(-1)$};
  % reflection arrows curving down between the lines
  \draw[->,trapred,dashed,thick] (1,-0.18) to[bend left=20] (-1,-1.62);
  \draw[->,trapred,dashed,thick] (3,-0.18) to[bend left=22] (-3,-1.62);
\end{tikzpicture}
\caption{Multiplying by $-1$ reflects the line through $0$. The blue point $a=1$
was left of the green $b=3$; after reflection $-a=-1$ is to the \emph{right} of
$-b=-3$. The order has flipped: $1<3$ but $-1>-3$.}
\label{ch40:fig-flip}
\end{figure}

\begin{trapbox}
The flip rule is the most common source of wrong inequality answers. When you
divide both sides of $-2x < 6$ by $-2$, you must reverse: $x > -3$, \emph{not}
$x<-3$. A safe habit: never multiply or divide an inequality by a quantity whose
sign you do not know. If you must divide by something like $x-1$, split into cases
$x>1$ and $x<1$ and handle each separately.
\end{trapbox}

\begin{example}[Solving a rational inequality the safe way]
\label{ch40:ex-rational-ineq}
Solve $\dfrac{x+1}{x-2} \ge 0$. Do \emph{not} cross-multiply --- you do not know
the sign of $x-2$. Instead find where each factor is zero or undefined: the
numerator vanishes at $x=-1$, the denominator at $x=2$ (where the expression is
undefined). These split the line into three regions. Test one point in each:
\[
  \begin{array}{c|ccc}
    \text{region} & x<-1 & -1<x<2 & x>2 \\\hline
    \text{sign of } \tfrac{x+1}{x-2} & (-)/(-)=+ & (+)/(-)=- & (+)/(+)=+
  \end{array}
\]
So the expression is $\ge 0$ on $x \le -1$ and $x > 2$ (including $x=-1$ where it
equals $0$, excluding $x=2$ where it is undefined). The solution is
$(-\infty,-1] \cup (2,\infty)$.
\end{example}

The whole method is a \emph{sign chart}: mark the points where the expression is
zero or undefined, then record the sign in each gap. Figure~\ref{ch40:fig-signchart}
shows it for this example.

\begin{figure}[htbp]
\centering
\begin{tikzpicture}[scale=1.0]
  \draw[->,thick] (-5,0)--(5,0) node[right]{$x$};
  % critical points
  \fill[targetgreen] (-1,0) circle (2.6pt);
  \node[below=3pt, targetgreen, font=\small] at (-1,0) {$x=-1$};
  \draw[fill=white,trapred,thick] (2,0) circle (2.6pt);
  \node[below=3pt, trapred, font=\small] at (2,0) {$x=2$};
  % zero / undefined annotations
  \node[above=3pt, targetgreen, font=\small] at (-1,0) {zero};
  \node[above=3pt, trapred, font=\small] at (2,0) {undefined};
  % signs in the three regions
  \node[sourceblue, font=\large] at (-3,0.5) {$+$};
  \node[trapred,    font=\large] at (0.5,0.5)    {$-$};
  \node[sourceblue, font=\large] at (3.6,0.5)  {$+$};
  % shaded solution set below with ample separation
  \draw[sourceblue, line width=2.6pt] (-5,-0.85)--(-1,-0.85);
  \fill[sourceblue] (-1,-0.85) circle (2.6pt);
  \draw[sourceblue, line width=2.6pt] (2,-0.85)--(5,-0.85);
  \draw[fill=white,sourceblue,very thick] (2,-0.85) circle (3.2pt);
  \node[sourceblue, font=\small, anchor=west] at (-5,-1.35)
       {solution of $\tfrac{x+1}{x-2}\ge 0$: $(-\infty,-1]\cup(2,\infty)$};
\end{tikzpicture}
\caption{Sign chart for $\tfrac{x+1}{x-2}$. The two critical points --- the zero at
$x=-1$ (filled, included) and the undefined point at $x=2$ (open, excluded) --- cut
the line into three regions. One test point per region fixes the sign there. The
heavy blue segments below the axis are the solution set, where the fraction is
$\ge 0$.}
\label{ch40:fig-signchart}
\end{figure}

\section{Absolute Value: Distance Made Algebraic}
\label{ch40:absolute-value}

The absolute value $\abs{x}$ is defined by cases,
\[
  \abs{x} = \begin{cases} x, & x \ge 0,\\ -x, & x < 0,\end{cases}
\]
but the case split hides the real idea, which is geometric:

\begin{intuitionbox}
$\abs{x}$ is the \emph{distance from $x$ to $0$} on the number line, and more
generally $\abs{x-a}$ is the \emph{distance between $x$ and $a$}. The two cases
say ``distance is never negative; if the number is already to the right of
$0$ its distance is itself, and if it is to the left we negate to make it
positive.'' Read every absolute value as a distance and most manipulations become
easy to see.
\end{intuitionbox}

This distance reading immediately explains how absolute-value inequalities unfold:
\[
  \abs{x-a} < r \quad\Longleftrightarrow\quad a - r < x < a + r,
\]
``$x$ is within $r$ of $a$,'' an \emph{interval} centered at $a$ with radius $r$.
And the opposite,
\[
  \abs{x-a} > r \quad\Longleftrightarrow\quad x < a-r \;\text{ or }\; x > a+r,
\]
``$x$ is farther than $r$ from $a$,'' the \emph{outside} of that interval.
Figure~\ref{ch40:fig-absval} draws the first case.

\begin{figure}[htbp]
\centering
\begin{tikzpicture}[scale=1.0]
  \draw[->,thick] (-1,0)--(9,0) node[right]{$x$};
  \foreach \x/\lab in {1/{a-r}, 4/{a}, 7/{a+r}}
     \draw (\x,0.09)--(\x,-0.09) node[below=2pt]{$\lab$};
  % shaded interval
  \draw[sourceblue, line width=3pt] (1,0)--(7,0);
  \fill[white,draw=sourceblue,thick] (1,0) circle (3pt);
  \fill[white,draw=sourceblue,thick] (7,0) circle (3pt);
  \fill[targetgreen] (4,0) circle (2.4pt);
  \node[above, sourceblue] at (4,0.35) {$\abs{x-a}<r$ : the open interval $(a-r,\,a+r)$};
  % radius braces
  \draw[decorate,decoration={brace,amplitude=5pt},quotientorange]
       (4,-0.7)--(1,-0.7) node[midway,below=4pt]{$r$};
  \draw[decorate,decoration={brace,amplitude=5pt},quotientorange]
       (7,-0.7)--(4,-0.7) node[midway,below=4pt]{$r$};
\end{tikzpicture}
\caption{$\abs{x-a}<r$ reads as ``$x$ is within distance $r$ of $a$,'' so the
solution set is the open interval $(a-r,a+r)$ --- a fixed center $a$ with a radius
$r$ in each direction. This exact picture is the geometry behind the
$\eps$--$\delta$ definition of a limit.}
\label{ch40:fig-absval}
\end{figure}

\begin{connectionbox}
Limits are defined with exactly these inequalities. The statement
$\lim_{x\to a} f(x) = L$ unpacks to: for every tolerance $\eps>0$ there is a
$\delta>0$ so that $\abs{x-a}<\delta$ forces $\abs{f(x)-L}<\eps$. Both halves are
absolute-value inequalities, i.e.\ ``within distance'' statements: \emph{whenever
$x$ is within $\delta$ of $a$, $f(x)$ is within $\eps$ of $L$.}
Figure~\ref{ch40:fig-absval}, drawn once for $x$ (radius $\delta$) and once for
$f(x)$ (radius $\eps$), is the picture behind the definition of a limit.
\end{connectionbox}

\begin{recoverybox}
\textbf{Pause and think:} If an error tolerance is written as $\abs{x - 5} < 0.02$,
what is the nominal target value, and what is the allowable range for $x$?
The nominal target is the center $a = 5$, and the maximum allowable error is the
radius $r = 0.02$. The allowable range is the open interval
$(5 - 0.02,\, 5 + 0.02) = (4.98,\, 5.02)$. Reading $\abs{x-a} < r$ as ``the distance
from $x$ to $a$ is less than $r$'' lets you write down the interval directly,
without working through the case split.
\end{recoverybox}

A few algebraic properties of the absolute value are used constantly in
\emph{analysis} (the careful, proof-based theory behind calculus).

\begin{focusbox}
\textbf{Useful identities.}
\[
  \abs{xy} = \abs{x}\,\abs{y}, \qquad \abs{x}^2 = x^2, \qquad \sqrt{x^2}=\abs{x}.
\]
\textbf{The triangle inequality.} For all real (indeed all complex) $x,y$,
\[
  \abs{x+y} \le \abs{x} + \abs{y}.
\]
\end{focusbox}

The triangle inequality turns up more than any other inequality in analysis.
Its name comes from the picture in the plane: the straight side of a triangle is no
longer than the sum of the other two sides. On the line it says that walking from
$0$ to $x$ and then on to $x+y$ is never shorter than walking straight from $0$ to
$x+y$. Figure~\ref{ch40:fig-triangle-ineq} draws the plane version.

\begin{whybox}
A proof of the triangle inequality for real numbers. For any real $t$, $-\abs t \le t
\le \abs t$ (a number lies between $\pm$ its size). Apply this to $x$ and to $y$:
\[
  -\abs x \le x \le \abs x, \qquad -\abs y \le y \le \abs y.
\]
Adding the two chains piece by piece gives
\[
  -(\abs x + \abs y) \le x + y \le \abs x + \abs y.
\]
A statement of the form $-c \le s \le c$ (with $c \ge 0$) is exactly
$\abs s \le c$. With $s = x+y$ and $c = \abs x + \abs y$ this reads
$\abs{x+y} \le \abs x + \abs y$. The proof rests on one observation:
every number lies between plus and minus its absolute value.
\end{whybox}

\begin{figure}[htbp]
\centering
\begin{tikzpicture}[scale=1.1]
  \coordinate (O) at (0,0);
  \coordinate (X) at (3.5,0.6);
  \coordinate (XY) at (4.8,2.8);

  % vectors / segments with lengths
  \draw[->,very thick,sourceblue] (O) -- (X) node[midway,below=3pt,font=\small] {$\abs{x}$};
  \draw[->,very thick,targetgreen!80!black] (X) -- (XY) node[midway,right=3pt,font=\small] {$\abs{y}$};
  \draw[->,very thick,quotientorange] (O) -- (XY) node[midway,above left=3pt,font=\small] {$\abs{x+y}$};

  % endpoints labeled by position
  \fill (O) circle (2pt) node[below left] {$0$};
  \fill (X) circle (2pt) node[below right,font=\small] {$x$};
  \fill (XY) circle (2pt) node[above right,font=\small] {$x+y$};

  % inequality statement callout
  \node[anchor=south,font=\small\bfseries,quotientorange] at (2.0,2.2) 
    {$\abs{x+y} \le \abs{x} + \abs{y}$};
\end{tikzpicture}
\caption{The triangle inequality in the plane: the straight path from $0$ to $x+y$, of length $\abs{x+y}$, is never longer than the detour through $x$, of length $\abs{x} + \abs{y}$. The two are equal exactly when $x$ and $y$ point in the same direction (for real numbers: have the same sign).}
\label{ch40:fig-triangle-ineq}
\end{figure}

\section{Functions: The Facts Calculus Leans On}
\label{ch40:functions}

Functions get a full chapter early in the book; here we collect, for reference, the
handful of facts the later chapters use without re-deriving.

\textbf{A function} $f\colon D \to \R$ assigns to each input $x$ in its domain $D$
\emph{exactly one} output $f(x)$. The ``exactly one'' is the defining demand --- it
is what the \emph{vertical line test} checks on a graph: a curve is the graph of a
function precisely when no vertical line meets it twice (which would mean a single
input generates multiple outputs).

\textbf{Composition.} $(g\circ f)(x) = g(f(x))$: do $f$ first, then feed the result
to $g$. Order matters; $g\circ f$ and $f\circ g$ are usually completely different
functions. The chain rule (Section~\ref{ch11:chain-rule}) is the
rule for differentiating a composition, so compositions appear in a large share
of all derivative calculations.

\textbf{Inverse functions.} $f^{-1}$ undoes $f$: $f^{-1}(f(x)) = x$ and
$f(f^{-1}(y)) = y$. An inverse exists exactly when $f$ is \emph{one-to-one} (each
output comes from only one input), which the \emph{horizontal line test} checks
(verifying that no horizontal line meets the graph twice, which would mean two
different inputs collide on the same output). The graph of $f^{-1}$ is the graph
of $f$ reflected across the diagonal line $y=x$, because inverting swaps the roles
of input and output, i.e.\ swaps the coordinates $(x,y) \mapsto (y,x)$, as
Figure~\ref{ch40:fig-inverse} shows.

\begin{figure}[htbp]
\centering
\begin{tikzpicture}
\begin{axis}[
    width=8.4cm, height=8.4cm,
    axis equal image,
    xlabel={$x$}, ylabel={$y$},
    xlabel style={at={(1,0)}, anchor=south west},
    ylabel style={at={(0,1)}, anchor=south west},
    domain=0:2.6, samples=120,
    xmin=-0.3, xmax=4.4, ymin=-0.3, ymax=4.4,
    axis lines=middle, grid=major, grid style={gray, opacity=0.2},
    legend style={at={(0.5,1.03)}, anchor=south, legend columns=-1, font=\footnotesize},
]
  \addplot[curveblue, very thick, domain=0:2] {x^2};
  \addlegendentry{$f(x)=x^2,\ x\ge0$}
  \addplot[targetgreen!75!black, very thick, domain=0:4] {sqrt(x)};
  \addlegendentry{$f^{-1}(x)=\sqrt{x}$}
  \addplot[gray, dashed, domain=-0.3:4.2] {x};
  \addlegendentry{$y=x$ (mirror)}
  \addplot[quotientorange, thick, dotted, forget plot] coordinates {(2,4) (4,2)};
  \addplot[only marks, mark=*, mark size=2pt, quotientorange, forget plot]
     coordinates {(2,4) (4,2)};
  \node[anchor=west, font=\small, quotientorange] at (axis cs:2.05,4.05) {$(2,4)$};
  \node[anchor=north, font=\small, quotientorange] at (axis cs:4,1.85) {$(4,2)$};
\end{axis}
\end{tikzpicture}
\caption{A function and its inverse are mirror images across the line $y=x$. Here
$f(x)=x^2$ (restricted to $x\ge0$ so it is one-to-one) and $f^{-1}(x)=\sqrt x$.
The orange points show one matched pair: $(2,4)$ is on the graph of $f$ because
$2^2=4$, and its mirror image $(4,2)$ is on the graph of $f^{-1}$ because
$\sqrt4=2$.}
\label{ch40:fig-inverse}
\end{figure}

\textbf{Even and odd symmetry.} $f$ is \emph{even} if $f(-x)=f(x)$ (graph
symmetric across the $y$-axis, like $x^2$ or $\cos x$); \emph{odd} if
$f(-x)=-f(x)$ (graph symmetric through the origin, like $x^3$ or $\sin x$). These
symmetries cut work in half: the integral of an odd function over a symmetric
interval $[-a,a]$ is automatically $0$.

\textbf{Monotonicity.} $f$ is \emph{strictly increasing} if $x_1<x_2$ always
implies $f(x_1)<f(x_2)$, and \emph{strictly decreasing} if it always implies
$f(x_1)>f(x_2)$. A function that is one or the other is called \emph{strictly
monotonic}. Such a function never takes the same value twice, so it is
one-to-one and has an inverse. The sign of the derivative tells you where a
function is increasing or decreasing, which links calculus back to this idea.

\begin{trapbox}
$f^{-1}$ means the \emph{inverse function}, not the reciprocal. $\sin^{-1}x$ is the
inverse sine (arcsine), not $1/\sin x$. The reciprocal is written $(\sin x)^{-1}$
or $1/\sin x = \csc x$. The clash of notation is unfortunate but standard, so
check which meaning is intended.
\end{trapbox}

\section{What a Proof Is, and Three Reliable Shapes}
\label{ch40:proof}

A theorem is a statement that has been proved; a proof is the argument that shows
it is true. Mathematics is stricter about this than any other subject: a proof is
a chain of
steps, each following from the previous by pure logic or from facts already
granted, leaving \emph{no} gap a skeptic could pry open. Checking a million cases
is not a proof if there is a million-and-first. In return, a proved theorem is
true in every case where its \emph{hypotheses} (the conditions it assumes) hold.

\begin{aestheticbox}
Why insist on proof when testing examples and spotting patterns works so often?
Because a pattern can hold for dozens of cases and still fail. Consider $f(n) = n^2 + n + 41$. For $n=0,1,2,\dots,39$ every
value is prime: forty cases in a row. But
$f(40)=40^2+40+41 = 41^2$ is not prime. Forty confirmations could not save a false
\emph{conjecture} (a statement believed but not yet proved). Calculus is built on
infinitely many cases (every $x$, every $\eps$), so we \emph{cannot} check them
one by one. Proof is the only tool that establishes a statement across infinitely
many cases at once.
\end{aestheticbox}

We will meet three shapes of proof again and again. Each is a strategy for getting
from hypotheses to conclusion.

\subsection{Direct proof}

The most natural shape: assume the hypothesis, and march forward by valid steps
until you reach the conclusion. Most ``if $P$ then $Q$'' theorems are proved this
way --- start from $P$, deduce, arrive at $Q$.

\begin{example}[A direct proof]
\label{ch40:ex-direct}
\emph{Claim: the sum of two even integers is even.} An even integer is one of the
form $2k$ for some integer $k$ (that is the definition we work from). So let the
two even integers be $m = 2j$ and $n = 2k$. Then
\[
  m + n = 2j + 2k = 2(j+k).
\]
Since $j+k$ is an integer, $m+n$ has the form $2 \times (\text{integer})$, which is
exactly the definition of even. Therefore $m+n$ is even. $\qed$

This is the standard structure of a direct proof: unfold the definitions in the
hypotheses, do careful algebra, and fold the result back into the definition used
in the conclusion.
\end{example}

\subsection{Proof by contradiction}

Sometimes the conclusion is awkward to reach head-on, but its \emph{negation} is
straightforward to analyze. So we assume the opposite of what we want, follow the
consequences, and reach a contradiction. Since valid reasoning from a true
premise can never produce a falsehood, our starting assumption (the negation) must
have been false --- so the original claim is true.

\begin{connectionbox}
The proof that $\sqrt 2$ is irrational is one of the earliest landmarks in deductive mathematics,
dating to the Pythagoreans in the fifth century BCE and recorded in Euclid's
\emph{Elements} (c.~300~BCE). The Pythagoreans held that
the geometric world was governed strictly by whole numbers and their ratios. The discovery that
the diagonal of a unit square could not be represented as any ratio of whole numbers ---
that it was \emph{incommensurable} (having no shared unit of measure with the side) ---
came as a shock. Legend holds that the discovery was kept secret under
penalty of death. Proof by contradiction forced
mathematicians to accept a fact that no measurement could ever have shown.
\end{connectionbox}

\begin{example}[$\sqrt 2$ is irrational --- proof by contradiction]
\label{ch40:ex-sqrt2}
\emph{Suppose, for contradiction, that $\sqrt 2$ is rational.} Then we can write
$\sqrt 2 = p/q$ as a fraction in \emph{lowest terms} --- meaning $p$ and $q$ share
no common factor (any common factor could be cancelled first). Square both sides:
\[
  2 = \frac{p^2}{q^2} \quad\Longrightarrow\quad p^2 = 2q^2.
\]
So $p^2$ is even, which forces $p$ itself to be even (if $p$ were odd, say $p = 2k+1$,
then $p^2 = (2k+1)^2 = 4k^2 + 4k + 1 = 2(2k^2+2k)+1$, which is odd). Write $p = 2r$.
Substituting,
\[
  (2r)^2 = 2q^2 \;\Longrightarrow\; 4r^2 = 2q^2 \;\Longrightarrow\; q^2 = 2r^2.
\]
Now $q^2$ is even, so $q$ is even too. But then $p$ and $q$ are \emph{both} even
--- they share the common factor $2$ --- contradicting our choice of a fraction in
lowest terms. The contradiction means our supposition was impossible: $\sqrt 2$
cannot be written as such a fraction. Hence $\sqrt 2$ is irrational. $\qed$
\end{example}

\begin{happeningbox}
Why were we allowed to assume lowest terms? Because cancelling common factors
makes the positive denominator strictly smaller, and positive integers cannot
shrink forever: every nonempty subset of $\N$ contains a least element (the
\emph{well-ordering principle}). So among all fractions equal to $p/q$ there is
one with the smallest denominator, and it must be in lowest terms. The same
principle is the foundation of mathematical induction, which comes next.
\end{happeningbox}

\subsection{Proof by induction}

Induction proves a statement $P(n)$ for \emph{every} natural number $n$ at once,
even though there are infinitely many. The image to hold is a line of dominoes.

\begin{intuitionbox}
Imagine infinitely many dominoes standing in a row, one for each $n=1,2,3,\dots$.
You want to know they all fall. You do not push each one --- there are infinitely
many. Instead you check two things: (1) the \textbf{first} domino falls, and
(2) \textbf{each} domino, if it falls, knocks over the next. Granted those two
facts, every domino must fall: the first falls by (1), which by (2) topples the
second, which topples the third, forever. That is induction.
\end{intuitionbox}

\begin{focusbox}
\textbf{Principle of mathematical induction.} To prove $P(n)$ holds for all
integers $n \ge 1$, prove two things:
\begin{itemize}
  \item \textbf{Base case:} $P(1)$ is true. (The first domino falls.)
  \item \textbf{Inductive step:} for every $k$, \emph{if} $P(k)$ is true \emph{then}
        $P(k+1)$ is true. (Each domino topples the next.)
\end{itemize}
The assumption ``$P(k)$ is true'' inside the step is called the \emph{inductive
hypothesis}. You are allowed to assume it precisely because you are about to show
it would force the next case.
\end{focusbox}

Figure~\ref{ch40:fig-domino} draws the two requirements on the dominoes.

\begin{figure}[htbp]
\centering
\begin{tikzpicture}[scale=0.95]
  % Ground line
  \draw[thick,gray!60] (-0.8,0) -- (10.8,0);

  % Falling domino 1
  \begin{scope}[shift={(0.6,0)},rotate around={-45:(0.45,0)}]
    \fill[sourceblue!25,rounded corners=2pt] (0,0) rectangle (0.45,1.8);
    \draw[thick,sourceblue,rounded corners=2pt] (0,0) rectangle (0.45,1.8);
    \node[black!85,font=\footnotesize\bfseries] at (0.225,0.9) {$1$};
  \end{scope}

  % Falling domino 2
  \begin{scope}[shift={(1.8,0)},rotate around={-35:(0.45,0)}]
    \fill[sourceblue!20,rounded corners=2pt] (0,0) rectangle (0.45,1.8);
    \draw[thick,sourceblue,rounded corners=2pt] (0,0) rectangle (0.45,1.8);
    \node[black!85,font=\footnotesize\bfseries] at (0.225,0.9) {$2$};
  \end{scope}

  % Dots
  \node[font=\bfseries,gray] at (3.2,0.8) {$\cdots$};

  % Domino k tilting
  \begin{scope}[shift={(4.7,0)},rotate around={-22:(0.45,0)}]
    \fill[quotientorange!25,rounded corners=2pt] (0,0) rectangle (0.45,1.8);
    \draw[thick,quotientorange,rounded corners=2pt] (0,0) rectangle (0.45,1.8);
    \node[black!85,font=\footnotesize\bfseries] at (0.225,0.9) {$k$};
  \end{scope}

  % Domino k+1 upright but about to be struck
  \begin{scope}[shift={(5.85,0)}]
    \fill[targetgreen!20,rounded corners=2pt] (0,0) rectangle (0.45,1.8);
    \draw[thick,targetgreen!80!black,rounded corners=2pt] (0,0) rectangle (0.45,1.8);
    \node[black!85,font=\footnotesize\bfseries] at (0.225,0.9) {$k{+}1$};
  \end{scope}

  % Domino k+2 upright
  \begin{scope}[shift={(7.3,0)}]
    \fill[gray!15,rounded corners=2pt] (0,0) rectangle (0.45,1.8);
    \draw[thick,black!50,rounded corners=2pt] (0,0) rectangle (0.45,1.8);
    \node[black!85,font=\footnotesize\bfseries] at (0.225,0.9) {$k{+}2$};
  \end{scope}

  % Dots
  \node[font=\bfseries,gray] at (8.7,0.8) {$\cdots$};

  % Annotations
  \draw[->,thick,sourceblue] (0.6,2.5) -- (1.0,1.5);
  \node[above,font=\small,sourceblue] at (0.6,2.5) {\textbf{Base case:} $P(1)$ true};

  \draw[->,thick,quotientorange] (5.9,2.5) -- (5.8,1.95);
  \node[above,font=\small,quotientorange] at (6.3,2.5) {\textbf{Inductive step:} $P(k) \Rightarrow P(k+1)$};
\end{tikzpicture}
\caption{Induction as falling dominoes. Proving $P(1)$ knocks over the first
domino. Proving $P(k) \Rightarrow P(k+1)$ means each falling domino knocks over
the next one. Together they make every domino fall, so $P(n)$ holds for all
integers $n \ge 1$.}
\label{ch40:fig-domino}
\end{figure}

\begin{example}[Sum of the first $n$ integers]
\label{ch40:ex-induction}
\emph{Claim:} for every $n \ge 1$,
\[
  1 + 2 + 3 + \cdots + n = \frac{n(n+1)}{2}.
\]
\textbf{Base case ($n=1$).} The left side is $1$; the right side is
$\tfrac{1\cdot 2}{2}=1$. They agree, so $P(1)$ holds.

\textbf{Inductive step.} Assume the formula holds for $n=k$ (the inductive
hypothesis):
\[
  1 + 2 + \cdots + k = \frac{k(k+1)}{2}.
\]
We must show it then holds for $n=k+1$. Add $(k+1)$ to both sides:
\[
  1 + 2 + \cdots + k + (k+1) = \frac{k(k+1)}{2} + (k+1).
\]
Factor $(k+1)$ out of the right-hand side:
\begin{align*}
  \frac{k(k+1)}{2} + (k+1) &= (k+1)\left(\frac{k}{2} + 1\right) \\
  &= (k+1)\cdot\frac{k+2}{2} = \frac{(k+1)(k+2)}{2}.
\end{align*}
That last expression is exactly the original formula with $n$ replaced by $k+1$.
So $P(k)\Rightarrow P(k+1)$. By induction the formula holds for all $n \ge 1$.
$\qed$
\end{example}

\begin{whybox}
A famous story has the young Carl Friedrich Gauss, assigned by his schoolmaster to sum
the integers from $1$ through $100$, producing the answer $5050$ almost instantly.
His insight was to pair the numbers from opposite ends:
\[
  1 + 100 = 101, \quad 2 + 99 = 101, \quad 3 + 98 = 101, \quad \dots, \quad 50 + 51 = 101.
\]
There are $50$ such pairs, each totaling $101$, giving $50 \times 101 = 5050$. For any
positive integer $n$, write the sum forward and backward:
\[
  \begin{array}{rcccccccc}
    S  &=& 1     &+& 2       &+& \cdots &+& n \\
    S  &=& n     &+& (n-1)   &+& \cdots &+& 1 \\ \hline
    2S &=& (n+1) &+& (n+1)   &+& \cdots &+& (n+1)
  \end{array}
\]
Adding the two lines produces $n$ columns, every one summing to $n+1$. Therefore
$2S = n(n+1)$, so $S = \tfrac{n(n+1)}{2}$. It helps to know a result two ways:
induction certifies that it holds without exception; Gauss's pairing explains \emph{why}
the formula takes that shape.
\end{whybox}

\begin{trapbox}
Induction without a base case proves nothing. ``If $P(k)$ then $P(k+1)$'' alone is
a chain of dominoes that nobody pushed --- the implications can all be valid while
the statement is false for every $n$. Always nail down the base case. Conversely, a
base case with a broken inductive step (a gap where $P(k)$ does not force
$P(k+1)$) is equally worthless: one fallen domino, then nothing.
\end{trapbox}

\begin{recoverybox}
\textbf{Pause and think:} Suppose someone claims: ``For all integers $n \ge 1$,
$n = n + 10$.'' They demonstrate that \emph{if} $k = k + 10$, then adding $1$ to
both sides yields $k + 1 = (k + 1) + 10$, so the inductive step holds. Why is the
claim false?
The inductive step is a valid conditional statement (``if true for $k$, then true
for $k+1$''), but the base case $P(1)$ is $1 = 11$, which is false. The
dominoes are spaced to knock each other over, but the first domino was never tipped.
\end{recoverybox}

\begin{connectionbox}
Induction is not confined to ``sum'' formulas. It underlies the power rule
$\frac{d}{dx}x^n = nx^{n-1}$ for positive integer $n$ (prove it for $n=1$, then
use the product rule to climb from $x^k$ to $x^{k+1}$), the binomial theorem, and
the convergence arguments behind Taylor series. Whenever a statement is indexed by
a whole number and each case rests on the one before, induction is the tool.
\end{connectionbox}

\section{The Binomial Theorem: Algebra and Counting Together}
\label{ch40:binomial}

One identity gets a section of its own because it appears throughout the
book: in the derivation of the power rule, in Taylor's theorem, and in
the definition of $e$. It answers: what is $(a+b)^n$ expanded?

We already know $(a+b)^2 = a^2 + 2ab + b^2$ and $(a+b)^3 = a^3 + 3a^2b + 3ab^2 +
b^3$. The coefficients $1,2,1$ and $1,3,3,1$ are rows of \emph{Pascal's
triangle}, where each entry is the sum of the two above it
(Figure~\ref{ch40:fig-pascal}).

\begin{figure}[htbp]
\centering
\begin{tikzpicture}[font=\small]
  \foreach \n/\row in {0/{1}, 1/{1,1}, 2/{1,2,1}, 3/{1,3,3,1}, 4/{1,4,6,4,1}} {
    \node[anchor=east] at (-3.4,-0.8*\n) {$(a+b)^{\n}$:};
    \foreach \v [count=\k from 0] in \row
      \node (p\n-\k) at ({(\k-\n/2)*1.1}, {-0.8*\n}) {$\v$};
  }
  \draw[->, thick, quotientorange, shorten <=1pt, shorten >=1pt] (p3-1) -- (p4-2);
  \draw[->, thick, quotientorange, shorten <=1pt, shorten >=1pt] (p3-2) -- (p4-2);
  \draw[quotientorange, thick] (p4-2) circle (0.25);
  \draw[->, thick, sourceblue, shorten <=1pt, shorten >=1pt] (p2-0) -- (p3-1);
  \draw[->, thick, sourceblue, shorten <=1pt, shorten >=1pt] (p2-1) -- (p3-1);
  \draw[sourceblue, thick] (p3-1) circle (0.25);
\end{tikzpicture}
\caption{Pascal's triangle. Row $n$ lists the coefficients of $(a+b)^n$. Each
entry is the sum of the two entries just above it: in blue $1+2=3$, in orange
$3+3=6$. The ends of every row are $1$.}
\label{ch40:fig-pascal}
\end{figure}

The general row is given by the binomial theorem. In it, $n!$ (``$n$
factorial'') is the product $1\cdot2\cdot3\cdots n$, with $0!=1$, and $\sum$
adds up the terms for $k=0,1,\dots,n$.

\begin{theorem}[Binomial theorem]
\label{ch40:thm-binomial}
For any nonnegative integer $n$,
\[
  (a+b)^n = \sum_{k=0}^{n} \binom{n}{k} a^{n-k} b^{k},
  \qquad \binom{n}{k} = \frac{n!}{k!\,(n-k)!}.
\]
\end{theorem}

\begin{happeningbox}
Where do the coefficients $\binom{n}{k}$ come from? Expanding $(a+b)^n =
(a+b)(a+b)\cdots(a+b)$ ($n$ factors) means choosing, from each factor, either the
$a$ or the $b$, then multiplying the choices. A term $a^{n-k}b^k$ arises every time
we pick $b$ from exactly $k$ of the $n$ factors. The number of ways to choose which
$k$ factors contribute a $b$ is $\binom{n}{k}$, read ``$n$ choose $k$'': the
number of ways to choose $k$ items from $n$. So the
binomial theorem states a counting fact in algebraic form: the
coefficient is the count of ways a given mix of $a$'s and $b$'s can occur. This is
why probability and algebra share the same triangle.
\end{happeningbox}

\begin{example}[Why it matters for the power rule]
\label{ch40:ex-binom-power}
The derivative of $x^n$ from the limit definition needs $(x+h)^n$. The binomial
theorem gives
\[
  (x+h)^n = x^n + n x^{n-1} h + \binom{n}{2} x^{n-2} h^2 + \cdots + h^n.
\]
Subtract $x^n$ and divide by $h$:
\[
  \frac{(x+h)^n - x^n}{h} = n x^{n-1} + \binom{n}{2}x^{n-2} h + \cdots + h^{n-1}.
\]
Every term after the first carries at least one factor of $h$, so as $h \to 0$ they
all vanish, leaving
\[
  \ddx{x^n} = n x^{n-1}.
\]
The binomial theorem isolated the one term that survives the limit: the term
with exactly one factor of $h$.
\end{example}

\section{Putting It Together: A Worked Gauntlet}
\label{ch40:gauntlet}

To see these tools cooperate, run one expression through several of them.

\begin{example}[Simplify a difference quotient]
\label{ch40:ex-gauntlet}
Simplify $\dfrac{1}{h}\!\left(\dfrac{1}{x+h} - \dfrac{1}{x}\right)$ and find its
limit as $h\to 0$ --- this is the derivative of $1/x$ from scratch.

\textbf{Step 1: combine the inner fractions} over the common denominator
$x(x+h)$:
\[
  \frac{1}{x+h} - \frac{1}{x} = \frac{x - (x+h)}{x(x+h)} = \frac{-h}{x(x+h)}.
\]
\textbf{Step 2: divide by $h$} (multiply by $1/h$); the $h$'s cancel, the move
that makes the limit tractable:
\[
  \frac{1}{h}\cdot\frac{-h}{x(x+h)} = \frac{-1}{x(x+h)}.
\]
\textbf{Step 3: take the limit} $h \to 0$. Now the expression is continuous at
$h=0$, so substitute:
\[
  \lim_{h\to 0} \frac{-1}{x(x+h)} = \frac{-1}{x \cdot x} = -\frac{1}{x^2}.
\]
So $\ddx{1/x} = -1/x^2$. Every step used only material from this appendix
(common denominators, care with signs, cancellation), yet the result is a
calculus fact.
\end{example}

\begin{labbox}
Run the same gauntlet on $\dfrac{\sqrt{x+h}-\sqrt{x}}{h}$. The right tool is the
conjugate: multiply top and bottom by $\sqrt{x+h}+\sqrt{x}$ and use the difference
of squares~\eqref{ch40:diff-squares} to turn the top into $(x+h)-x = h$, which
cancels the $h$ below. You should land on $\dfrac{1}{2\sqrt x}$ in the limit ---
the derivative of $\sqrt x$. It is the same three-step procedure: rewrite,
cancel the $h$, substitute.
\end{labbox}

\section{The Scope and Limits of Finite Algebra}
\label{ch40:what-next}

The tools in this appendix turn intuitive steps into checkable ones. Factoring,
multiplying by conjugates, and combining fractions let us rewrite a difference
quotient $\frac{f(x+h)-f(x)}{h}$ so that the factor $h$ cancels before we take
the limit. Extending exponents to zero, negative, and fractional powers brings
reciprocals and roots under the single power rule $\ddx{x^r} = r x^{r-1}$.
Inequalities and absolute values state tolerances precisely: ``$\abs{x - a} <
\delta$ forces $\abs{f(x) - L} < \eps$'' is the definition of a limit. Induction
and contradiction let us prove claims about infinitely many cases at once.

But algebra alone cannot carry calculus. Addition, multiplication, factoring, and
expanding are finite operations, and they work just as well if the only numbers
we allow are the rationals $\Q$: every field axiom and every identity in this
appendix still holds there. Calculus, however, fails in $\Q$.

The rational line has gaps. As the contradiction proof of
Example~\ref{ch40:ex-sqrt2} showed, $\sqrt{2}$ is not rational. The rationals
$1, 1.4, 1.41, 1.414, \dots$ get closer and closer together, yet there is no
rational number they approach. In $\Q$ the Intermediate Value Theorem fails, as
Figure~\ref{ch40:fig-rational-gap} shows: $f(x) = x^2 - 2$ goes from $f(1) = -1$
to $f(2) = 2$ without ever equaling zero, because the only point in between where
$x^2-2=0$ is $\sqrt{2}$, which is missing from $\Q$. Maximum values can go missing
too: on the rationals between $1$ and $2$, the function $g(x) = -(x^2-2)^2$ gets
as close to $0$ as we like but never reaches it, because the \emph{supremum}
(least upper bound) of its values, $0$, would be reached only at $x=\sqrt2$.

\begin{figure}[htbp]
\centering
\begin{tikzpicture}
\begin{axis}[
    width=10cm, height=6cm,
    axis lines=middle, xmin=0.6, xmax=2.3, ymin=-1.5, ymax=2.6,
    xtick={1,2}, ytick={-1,1,2}, clip=false,
    xlabel={$x$}, ylabel={$y$},
    xlabel style={anchor=west}, ylabel style={anchor=south},
    tick label style={font=\small},
    domain=0.8:2.1, samples=100,
]
  \addplot[curveblue, very thick] {x^2-2};
  \addplot[only marks, mark=*, mark size=2pt, sourceblue]
     coordinates {(1,-1) (2,2)};
  \node[anchor=north west, font=\small] at (axis cs:1.02,-1.02) {$f(1)=-1$};
  \node[anchor=east, font=\small] at (axis cs:1.97,2) {$f(2)=2$};
  \draw[trapred, thick, fill=white] (axis cs:1.41421,0) circle (3pt);
  \draw[->, trapred] (axis cs:1.6,-0.9) -- (axis cs:1.44,-0.1);
  \node[anchor=north west, font=\small, trapred, align=left]
     at (axis cs:1.55,-0.85) {$x=\sqrt2$:\\ missing from $\Q$};
\end{axis}
\end{tikzpicture}
\caption{The graph of $f(x)=x^2-2$ runs from below the $x$-axis to above it.
Over the real numbers it crosses at $x=\sqrt2$. If only rational inputs are
allowed, that point (open red circle) is absent, so the curve passes from
negative to positive without ever taking the value $0$.}
\label{ch40:fig-rational-gap}
\end{figure}

\begin{connectionbox}
To make calculus work, we need the \emph{completeness} of the real numbers. That
is the subject of the next chapter, \emph{Appendix: The Real Numbers and
Completeness}. It shows how the gaps in $\Q$ are filled to build $\R$, and states
the \emph{least upper bound property}: every nonempty set of real numbers that
has an upper bound has a smallest upper bound. That property guarantees that
every increasing sequence with an upper bound converges, that a continuous curve
going from negative to positive crosses zero, and that every definite integral of
a continuous function on a closed interval is a real number.
\end{connectionbox}

\section*{Exercises}
\addcontentsline{toc}{section}{Exercises}

\begin{exercise}
\label{ch40:exr-factor}
Factor completely over the reals:
(a) $x^3 - 4x$;\quad (b) $x^4 - 16$;\quad (c) $x^2 - 6x + 9$;\quad
(d) $8x^3 - 27$. \\
For each, state how you recognized the pattern (common factor, difference of
squares, perfect square, difference of cubes).
\end{exercise}

\begin{exercise}
\label{ch40:exr-exponents}
Simplify to a single power (no negative or fractional exponents in the final
answer where avoidable):
(a) $\dfrac{a^5 a^{-2}}{a^{-4}}$;\quad
(b) $\bigl(8x^6\bigr)^{1/3}$;\quad
(c) $\dfrac{(2x^2)^3}{4x^4}$;\quad
(d) $\dfrac{x^{1/2}\,x^{1/3}}{x^{-1/6}}$.
Explain at each step which exponent law you are invoking.
\end{exercise}

\begin{exercise}
\label{ch40:exr-quadratic}
Solve $2x^2 - 4x - 3 = 0$ two ways: once with the quadratic formula, and once by
completing the square. Confirm the answers agree, and state the discriminant and
what it tells you about the roots.
\end{exercise}

\begin{exercise}
\label{ch40:exr-inequality}
Solve each inequality, writing the answer in interval notation, and say where (if
anywhere) the ``flip'' rule was needed:
(a) $3 - 2x \le 11$;\quad
(b) $\dfrac{x-1}{x+2} \le 0$;\quad
(c) $x^2 < 5x - 6$.
\end{exercise}

\begin{exercise}
\label{ch40:exr-absval}
Rewrite each as an interval or union of intervals on the number line, and sketch
it:
(a) $\abs{x-3} < 2$;\quad
(b) $\abs{2x+1} \ge 5$;\quad
(c) $\abs{x} = \abs{x-4}$ (interpret this one as ``equidistant from $0$ and $4$''
before computing).
\end{exercise}

\begin{exercise}
\label{ch40:exr-triangle}
Use the triangle inequality to prove the \emph{reverse triangle inequality}
$\bigl|\,\abs{x}-\abs{y}\,\bigr| \le \abs{x-y}$.
(Hint: write $x = (x-y) + y$ and apply $\abs{x+y}\le\abs{x}+\abs{y}$ to get
$\abs{x} \le \abs{x-y} + \abs{y}$; rearrange, then do the same with the roles of
$x$ and $y$ swapped.)
\end{exercise}

\begin{exercise}
\label{ch40:exr-inverse}
Let $f(x) = \dfrac{2x-1}{x+3}$.
(a) Find $f^{-1}$ by solving $y = f(x)$ for $x$.
(b) Verify $f(f^{-1}(y)) = y$.
(c) State the domain of $f$ and the domain of $f^{-1}$, and explain why they are
what they are.
\end{exercise}

\begin{exercise}
\label{ch40:exr-direct}
Give a direct proof that the product of two odd integers is odd. (Start from the
definition: an odd integer has the form $2k+1$.)
\end{exercise}

\begin{exercise}
\label{ch40:exr-contradiction}
Prove by contradiction that $\sqrt 3$ is irrational. Identify the one place where
the argument differs from the $\sqrt 2$ proof of Example~\ref{ch40:ex-sqrt2}
(hint: you will use that ``$3 \mid p^2 \Rightarrow 3 \mid p$,'' where $3 \mid p$
means ``$3$ divides $p$ with no remainder''; this step deserves a sentence of
justification).
\end{exercise}

\begin{exercise}
\label{ch40:exr-induction-squares}
Prove by induction that for all $n \ge 1$,
\[
  1^2 + 2^2 + 3^2 + \cdots + n^2 = \frac{n(n+1)(2n+1)}{6}.
\]
Be explicit about your base case, your inductive hypothesis, and the algebra that
turns the $n=k$ formula into the $n=k+1$ formula.
\end{exercise}

\begin{exercise}
\label{ch40:exr-induction-ineq}
Prove by induction that $2^n > n$ for every integer $n \ge 1$. (This is an
\emph{inequality} induction; in the step, add to or multiply the inductive
hypothesis to reach the $k+1$ case.)
\end{exercise}

\begin{exercise}
\label{ch40:exr-binomial}
Use the binomial theorem (Theorem~\ref{ch40:thm-binomial}) to expand $(1+x)^4$,
and verify your coefficients against the corresponding row of Pascal's triangle.
Then use the same expansion to write down the first three terms of the difference
quotient $\bigl((1+h)^4 - 1\bigr)/h$ and read off the derivative of $x^4$ at
$x=1$.
\end{exercise}
